\documentclass[10pt,a4paper]{amsart}
\usepackage[margin=19mm]{geometry}
\usepackage{amsmath,amssymb,mathtools}
\usepackage[scr=rsfs,cal=euler]{mathalfa}
\usepackage{xfrac}
\usepackage[unicode,hidelinks,bookmarks=false]{hyperref}
\newcommand{\ie}{i.e.\ }
\newcommand{\cf}{cf.\ }
\newcommand{\resp}{respectively\ }
\newcommand{\Subsection}{\subsection}
\providecommand{\nobreakdash}{\nobreak\hbox{-}}
\setlength{\emergencystretch}{2em}
\multlinegap0pt
\usepackage[shortlabels]{enumitem}
\usepackage{amssymb}
\usepackage{mathtools}
\usepackage{float}
\newtheorem{theorem}{Theorem}
\title{Dynamic-memory fractional calculus via generator-based memory construction: operational theory, semigroup structure, and applications}
\author{Jehad Alzabut}
\date{Source: arXiv:2605.26326v1 (CC BY 4.0)}
\begin{document}
\maketitle

\noindent\emph{Source and licence.} Dynamic-memory fractional calculus via generator-based memory construction: operational theory, semigroup structure, and applications. Version arXiv:2605.26326v1 (CC BY 4.0), distributed by arXiv under the Creative Commons Attribution 4.0 International licence, http://creativecommons.org/licenses/by/4.0/, as recorded in the arXiv licence metadata for this exact version and shown on the abstract page https://arxiv.org/abs/2605.26326v1. Full licence notice: \texttt{licenses/arxiv-2605.26326-CC-BY-4.0.txt}.\par\smallskip

\noindent\textit{Editorial scope.} Counted unit: the article's theorem theorem (source0.tex lines 1722--1758). The article's complete original proof of it is delivered with it (source0.tex lines 1760--1817, one \texttt{proof} environment of 58 lines). No mathematics is written, repaired or re-derived: the statement and the proof are the article's own lines, in the article's own notation, and no retained line is substituted or re-flowed.  No further source-local context is needed: every cross-reference of every retained line resolves inside the two delivered spans.  Nothing was omitted from any retained span, so the package carries no omission record.  No retained line contains a citation, so the wrapper carries no bibliography and no bibliography entry.  No retained line contains a figure environment, an image-inclusion command or drawing code, so no image is delivered and the package declares no asset dependency.  Editorial formatting only, in addition: an AMS \texttt{amsart} preamble that restates the article's own declarations and macros used by the retained lines, namely the environments \texttt{theorem} on separate counters, together with the standard packages those lines need.  The retained environments print in this document as Theorem 1 in order of appearance, whereas the article prints the same items with the numbering of its own full document.  The complete native source of the article as distributed in the arXiv e-print for this exact version is bundled unchanged as \texttt{sources/201408/source0.tex}.
\begin{theorem}[Fundamental properties of the $\mathcal{DMFI}$ and $\mathcal{DMFD}$]
Suppose that the semigroup property holds and that the involved operators are well defined.
Then the following properties hold:
\begin{enumerate}[label=(\roman*),itemsep=2pt,topsep=2pt]
\item
\(
\mathcal D_\Phi^\alpha \mathcal I_\Phi^\alpha x=x.
\)

\item
For sufficiently smooth functions,
\(
\mathcal I_\Phi^\alpha\,{}^c \mathcal D_\Phi^\alpha x=x-x(0).
\)

\item
The operator \(\mathcal I_\Phi^\alpha\) is linear.

\item
If \(\Psi_\alpha(\cdot;\Theta)\in L^1(0,\mathcal{T})\), then
\(
\mathcal I_\Phi^\alpha:L^\infty(0,\mathcal{T})\to L^\infty(0,\mathcal{T})
\)
is bounded.

\item
Assume that the kernels \(\Psi_{1-\alpha}(\cdot;\Theta)\) form an approximate identity as
\(\alpha\to1^{-}\), that is,
\[
\int_0^t \Psi_{1-\alpha}(t-s;\Theta)f(s)\,ds \longrightarrow f(t)
\]
for every continuous function \(f\) on \([0,\mathcal{T}]\). Then, for every \(x\in \mathcal{C}^1([0,\mathcal{T}])\),
\[
\lim_{\alpha\to1^-}{}^c \mathcal{D}_\Phi^\alpha x(t)=x'(t).
\]
\end{enumerate}
\end{theorem}

\begin{proof}
Property (i) follows from the inverse relation between the $\mathcal{DMFI}$ and the corresponding $\mathcal{RL\;DMFD}$, together
with the semigroup property.

For (ii), by the definition of the $\mathcal{C\;DMFD}$,
\[
{}^c \mathcal D_\Phi^\alpha x=\mathcal I_\Phi^{1-\alpha}x'.
\]
Applying \( \mathcal I_\Phi^\alpha\) and using the semigroup property gives
\[
\mathcal I_\Phi^\alpha\,{}^c \mathcal D_\Phi^\alpha x
=
\mathcal I_\Phi^\alpha \mathcal I_\Phi^{1-\alpha}x'
=
\mathcal I_\Phi^1x'.
\]
Hence,
\[
\mathcal I_\Phi^\alpha\,{}^cD_\Phi^\alpha x
=
\int_0^t x'(s)\,ds
=
x(t)-x(0).
\]
Property (iii) follows directly from the linearity of the convolution integral defining
\(\mathcal I_\Phi^\alpha\).

To prove (iv), let \(x\in L^\infty(0,\mathcal T)\). Then
\[
|(\mathcal I_\Phi^\alpha x)(t)|
\le
\int_0^t |\Psi_\alpha(t-s;\Theta)|\,|x(s)|\,ds
\le
\|x\|_\infty \int_0^t |\Psi_\alpha(t-s;\Theta)|\,ds.
\]
Therefore,
\[
\|\mathcal I_\Phi^\alpha x\|_\infty
\le
\|\Psi_\alpha\|_{L^1(0,\mathcal T)}\|x\|_\infty,
\]
which proves the boundedness.

Finally, for (v), from the definition of the $\mathcal{C\;DMFD}$ we have
\[
{}^c \mathcal D_\Phi^\alpha x(t)
=
\int_0^t \Psi_{1-\alpha}(t-s;\Theta)x'(s)\,ds.
\]
Since \(x\in \mathcal C^1([0,\mathcal T])\), the function \(x'\) is continuous on \([0,\mathcal T]\). Hence, by the
approximate identity assumption applied to \(f=x'\), we obtain
\[
\lim_{\alpha\to1^-}{}^c \mathcal D_\Phi^\alpha x(t)
=
x'(t).
\]
This completes the proof.
\end{proof}

\end{document}
